\section{Experimental Setup}
\label{sec:experimental_setup}

\subsection{Data}
\label{sec:exp_datasets}
We use the warehouse scenes of the NVIDIA PhysicalAI-SmartSpaces benchmark (MTMC tracking data)~\cite{wang2026smartspaces}, rendered with NVIDIA Omniverse, and partition them by facility. Training uses two cameras of \texttt{Warehouse\_000} sampled at $0.5$\,FPS ($300$ frames, $12{,}959$ boxes); validation uses two cameras of the different facility \texttt{Warehouse\_015} ($100$ frames, $3{,}696$ boxes). Sampling at $0.5$\,FPS reduces the redundancy between consecutive frames and keeps CPU training tractable; with static cameras, frames $2$\,s apart remain strongly correlated (Section~\ref{sec:limitations}). No facility, camera, or video sequence is shared, but 3D assets such as forklift and robot models come from the same library, so object appearance is not independent across partitions, and all views are static simulated cameras. The training/validation instances are \texttt{worker} $11{,}399$/$3{,}488$, \texttt{forklift} $600$/$150$, \texttt{amr} $724$/$45$, and \texttt{pallet\_truck} $236$/$13$; the validation partition is therefore strongly imbalanced ($94.37\%$ workers, $0.35\%$ pallet trucks). Boxes with a side of at most $4$ pixels are removed and frames are letterboxed from $1920 \times 1080$ to $640 \times 640$. All arms use the Ultralytics default augmentation, in which mosaic is switched off for the last $10$ epochs (\texttt{close\_mosaic=10}).

\subsection{Experimental Arms and Training}
\label{sec:exp_baselines}
\textbf{Arm~1} is the unmodified YOLO11s ($9.43$\,M parameters, $21.7$ GFLOPs). \textbf{Arm~2} trains SS2D-P5 ($8.80$\,M, $21.2$ GFLOPs) in a single stage, \textbf{Arm~3} trains SS2D-P5 with the two-stage protocol, and \textbf{Arm~4} trains SS2D-P4P5 ($8.89$\,M, $21.8$ GFLOPs) in a single stage. All arms are fine-tuned from the same COCO checkpoint for $20$ epochs; only the SS2D blocks start untrained. The comparison is not a single-variable intervention: \texttt{C2PSA} starts pretrained whereas the SS2D blocks receive only about $1{,}500$ optimizer steps, and Arm~3 differs from Arm~2 both in freezing and in its restarted schedule.

Training used Python 3.14.8, PyTorch 2.14.1, and Ultralytics (8.4.185 for Arms 1--3 and 8.4.187 for Arm~4, read from the checkpoints), with AdamW selected by \texttt{optimizer=auto} ($\eta = 0.00125$, $\beta_1 = 0.9$, linear decay to $0.01\eta$ after a $3$-epoch warmup, weight decay $0.0005$), batch size $4$, input size $640$, and seed $0$ (one run per arm) on a host CPU workstation. We report the final-epoch model and the checkpoint with the highest validation fitness ($0.1\,\text{mAP}_{50} + 0.9\,\text{mAP}_{50-95}$); because both are evaluated on the partition used for selection, best-fitness figures are optimistically biased. Metrics are class-macro $\text{mAP}_{50}$ and $\text{mAP}_{50-95}$, and precision and recall at the confidence threshold that maximizes F1.

\subsection{Edge Measurement Protocol}
\label{sec:exp_export}
Checkpoints are exported to ONNX (opset 17, \texttt{onnxslim}) with standard operators only. The Pi files are the exports of the trained best checkpoints of Arm~1 (baseline), Arm~3 (SS2D-P5), and Arm~4 (SS2D-P4P5), identified by SHA-256 hashes; an earlier SS2D-P5 measurement that mistakenly used a pre-fine-tuning export is superseded. Models run in ONNX Runtime 1.31 (CPU execution provider, FP32, \texttt{ORT\_ENABLE\_ALL}, batch size $1$) on an Intel Core i7-7600U host (default threads) and on a Raspberry Pi~4 Model~B Rev~1.2 (Cortex-A72, $1.5$\,GHz, 4\,GB; Debian 13, Python 3.13.5; no heatsink or fan; official 5\,V/3\,A power supply) with $1$ or $4$ intra-op threads. We call a model \emph{edge-feasible} on the Pi if it runs there in FP32 without custom operators or plugins and within $1$\,GiB of peak memory; this criterion does not imply real-time throughput.

\textbf{Latency} is the time of \texttt{session.run} on a fixed input after untimed warmup, without preprocessing or NMS; we report the median, mean $\pm$ SD, and FPS $= 1000/\text{median}$. Single-thread runs use $10$ warmup and $100$ timed calls; four-thread runs use $5$ warmup and $30$ timed calls, because $100$-call runs heated the SoC to $79$--$84\,^{\circ}$C and once lowered the clock to $1.23$--$1.48$\,GHz; those runs are discarded. Every reported run started below $55\,^{\circ}$C. The Pi results come from three sessions: session~1 (baseline and SS2D-P4P5; clock logged only for the four-thread runs, which stayed at $1.5$\,GHz), session~2 (baseline, SS2D-P5, and the untrained intervention below; clock logged throughout, $1.5$\,GHz whenever a model ran, at most $76.9\,^{\circ}$C), and session~3 (the parallel scan below with its controls; clock logged throughout, at most $75.9\,^{\circ}$C). Relative values use the baseline of the same session. \textbf{Memory} is the peak resident set size of the process (\texttt{VmHWM}, including Python and session creation) on the Pi, and the RSS after the timed loop on the host. \textbf{Operator profiles} use the ONNX Runtime profiler with four threads and $5$ profiled calls after $2$ warmup calls.

\textbf{Runtime intervention.} To test whether graph size drives latency, we export the same SS2D block in three graph forms: the unrolled sequential scan used above, the same recurrence as one ONNX \texttt{Loop} per direction, and a chunked scan with chunk size $T = 4$. As runtime references we also export the baseline, the model with the $P_5$ block removed (identity), and a convolution-only replacement. These exports use untrained weights (seed $0$, BN statistics calibrated on random inputs, $\gamma = 1$), because they are compared only with each other. The three SS2D forms produce the same outputs (maximum relative difference $2.3 \times 10^{-6}$, versus $0.154$ between the SS2D and convolution-only exports as a negative control). Each runs with $1$ and $4$ threads, $3$ warmup and $20$ timed calls, in session~2. To test whether the sequential dependency chain drives latency, we also export the trained Arm~3 weights with the scan evaluated as a parallel (Hillis--Steele) associative scan~\cite{blelloch1990prefix}, which combines each position with the one $k = 1, 2, 4, \dots$ steps earlier and thus needs $\lceil \log_2 L \rceil$ steps instead of $L$; it matches the sequential recurrence to $2.8 \times 10^{-14}$ in float64 and the measured Arm~3 export to a relative output difference of $1.9 \times 10^{-7}$. It runs in session~3 with the trained baseline and unrolled Arm~3 exports as controls, under the protocol of the main measurements.
